\begin{musica}{A Banda}{Chico Buarque}{}
\nota{Refrão:}
Es\ac{D6/9}{ta}va à toa na \ac{A7}{vi}da, o meu amor me cha\ac{F#m7}{mou}
P’ra ver a \ac{B7}{ban}da pas\ac{E7/9}{sar}, cantando \ac{A7}{coi}sas de a\ac{D6/9}{mor};
A minha gente so\ac{A7}{fri}da, despediu-se da \ac{F#m7}{dor}
P’ra ver a \ac{B7}{ban}da pas\ac{E7/9}{sar}, cantando \ac{A7}{coi}sas de a\ac{D6/9}{mor}.

O homem \ac{Dmaj7}{sé}rio que contava \ac{A7}{di}nheiro, parou
O faro\ac{Am6/C}{lei}ro que \ac{B7}{con}tava \ac{Em7}{van}tagens, pa\ac{Em/D}{rou}
A na\ac{C#m7}{mo}rada que con\ac{F#7(9)}{ta}va as es\ac{F#m7}{tre}las,
Pa\ac{B7}{rou} para \ac{E7/9}{ver}, ouvir e dar \ac{Em7(9)}{pa}ssa\ac{A7}{gem}.

A moça \ac{Dmaj7}{tris}te que vivia \ac{A7}{ca}lada, sorriu
A rosa \ac{Am6/C}{tris}te que \ac{B7}{vi}via fe\ac{Em7}{cha}da, se a\ac{Em/D}{briu}
A meni\ac{C#m7}{na}da toda se assa\ac{F#7(9)}{nhou}
\ac{F#m7}{P’ra} ver a \ac{B7}{ban}da pas\ac{E7/9}{sar}, cantando \ac{A7}{coi}sas de a\ac{D}{mor}.

\nota{Refrão}

O velho fraco se esqueceu do cansaço e pensou
Qu'inda era moço pra sair no terraço e dançou
A moça feia debruçou na janela
Pensando que a banda tocava pra ela

A marcha alegre se espalhou na avenida e insistiu
A lua cheia que vivia escondida, surgiu
Minha cidade toda se enfeitou
Pra ver a banda passar, cantando coisas de amor

\ac{D6/9}{Mas} para meu de\ac{A7}{sen}canto, o que era doce aca\ac{F#m7}{bou}
Tudo \ac{B7}{to}mou seu lu\ac{E7/9}{gar}, depois que a \ac{A7}{ban}da pas\ac{D6/9}{sou}
E cada qual no seu \ac{A7}{can}to, em cada canto uma \ac{F#m7}{dor}
Depois da \ac{B7}{ban}da pas\ac{E7/9}{sar}, cantando \ac{A7}{coi}sas de a\ac{D6/9}{mor},
Depois da \ac{B7}{ban}da pas\ac{E7/9}{sar}, cantando \ac{A7}{coi}sas de a\ac{D6/9}{mor},
Depois da \ac{B7}{ban}da pas\ac{E7/9}{sar}, cantando \ac{A7}{coi}sas de a\ac{D6/9}{mor}.

\posicoes{C#m6: x46454, D6/9: x5445x, Dmaj7: xx0222, E7/9: xx2132, Em7: 022030,
  Em7(9): x7577x, Em/D: xx0453, F#7(9): xx4354, F#m7: 242222, A7: x02020,
  Am6/C: x34252, B7: x21202}
\end{musica}
